\section*{Bab \ref{ch:b1:observing-the-sky} --- Teleskop, detektor dan fotometri}

\begin{solution}{exo:b1:observing-the-sky:1}
Selisih 5 magnitudo sama dengan faktor 100: bintang bermagnitudo 1.0 mengirimkan 100 kali lebih banyak \omterm{def:b1:observing-the-sky:flux}{fluks}. Rigel dan Polaris: $10^{0.4\times(2.02 - 0.13)} =
10^{0.756} = 5.70$; Rigel, yang magnitudonya lebih kecil, lebih terang.
\end{solution}

\begin{solution}{exo:b1:observing-the-sky:2}
$d = 1/0.19495 = \qty{5.13}{\parsec}$, yaitu $5.1295\times206265 =
\qty{1.058e6}{\au}$ dan $5.1295\times3.2616 = \qty{16.7}{\lightyear}$.
\end{solution}

\begin{solution}{exo:b1:observing-the-sky:3}
$1.22\times\num{550e-9}/2.4 = \qty{2.80e-7}{rad} = \ang{;;0.058}$, sekitar
tujuh kali lebih halus daripada \omterm{def:b1:observing-the-sky:seeing}{kekaburan} \ang{;;0.40}: di angkasa teleskop
mencapai batas difraksinya. Untuk \qty{6.5}{m} pada \qty{2.0}{\micro\metre}:
$1.22\times\num{2.0e-6}/6.5 = \qty{3.75e-7}{rad} = \ang{;;0.077}$.
\end{solution}

\begin{solution}{exo:b1:observing-the-sky:4}
$B - V$ naik dari Rigel ($-0.03$) ke Altair (0.22), Polaris (0.60) dan Betelgeuse (1.85): yang paling panas adalah Rigel, lalu Altair, Polaris dan Betelgeuse. Jarak meredupkan kedua pita dengan faktor yang sama, sehingga $B - V$ tidak berubah olehnya (selama tidak ada debu yang memerahkan cahaya di perjalanan).
\end{solution}

\begin{solution}{exo:b1:observing-the-sky:5}
$X_1 = \sec\ang{30} = 1.1547$, $X_2 = \sec\ang{60} = 2.0000$.
$k = (6.590 - 6.412)/(2.0000 - 1.1547) = 0.178/0.8453 = \num{0.211}$ per
\omterm{def:b1:observing-the-sky:extinction}{massa udara}, dan $m_0 = 6.412 - 0.2106\times1.1547 = 6.169$.
\end{solution}

\begin{solution}{exo:b1:observing-the-sky:6}
$\mathrm{SNR} = \sqrt{400} = 20$. SNR tumbuh sebagai $\sqrt t$: untuk 100,
$t = \qty{10}{s}\times(100/20)^2 = \qty{250}{s}$. Galat \num{0.01}
magnitudo memerlukan $\mathrm{SNR} = 1.086/0.01 = 108.6$, sehingga
$t = 10\times(108.6/20)^2 = \qty{295}{s}$.
\end{solution}

\begin{solution}{exo:b1:observing-the-sky:7}
Jaraknya adalah $d = 1/0.00378 = \qty{264.6}{\parsec}$, dan \omterm{def:b1:observing-the-sky:absolute}{magnitudo mutlaknya} $M = 0.13 - 5\log_{10}26.46 = 0.13 - 7.11 = -6.98$. $L/L_\odot = 10^{0.4\times(4.83 + 6.98)} =
10^{4.724} = \num{5.3e4}$ dalam V.
\end{solution}

\begin{solution}{exo:b1:observing-the-sky:8}
$\pi\times1.2^2 = \qty{4.52}{m^2}$, dan $25/4.52 = 5.5$: teleskop Webb mengumpulkan jumlah foton yang sama dalam pemaparan 5.5 kali lebih pendek. Cakram penuh berdiameter \qty{6.5}{m} akan memiliki $\pi\times3.25^2 = \qty{33.2}{m^2}$; cermin sebenarnya memiliki $25/33.2 = \qty{75}{\percent}$ darinya, karena 18 heksagon meninggalkan tepi cakram itu tak tertutup dan celah antarsegmen serta bukaan di tengah tidak mengumpulkan apa pun.
\end{solution}

\begin{solution}{exo:b1:observing-the-sky:9}
(a) $\qty{10}{\parsec} = \qty{3.0857e17}{m}$, dan $F = \num{3.828e26}/(4\pi
\times(\num{3.0857e17})^2) = \num{3.828e26}/\num{1.1965e36} =
\qty{3.20e-10}{W.m^{-2}}$, $1361/\num{3.20e-10} = \num{4.3e12}$ kali lebih lemah daripada di Bumi.
(b) $5\log_{10}(d/\qty{10}{\parsec}) = -26.74 - 4.83 = -31.57$, sehingga
$d = \qty{10}{\parsec}\times10^{-6.314} = \qty{4.853e-6}{\parsec}$, yaitu
$\num{4.853e-6}\times206265 = \qty{1.001}{\au}$.
\end{solution}

\begin{solution}{exo:b1:observing-the-sky:10}
Di atas atmosfer \omterm{def:b1:observing-the-sky:flux}{fluksnya} $F_0$; di bawahnya, $F = F_0e^{-\tau_0X}$. Maka
$m - m_0 = -2.5\log_{10}(F/F_0) = 2.5\log_{10}e\;\tau_0X$, yang sama dengan $kX$ dengan
$k = 2.5\log_{10}e\;\tau_0 = 1.0857\,\tau_0$. Untuk $k = 0.160$:
$\tau_0 = 0.160/1.0857 = 0.147$. Di \omterm{def:b1:spherical-astronomy:celestial-sphere}{zenit} $e^{-0.147} = 0.863$ cahaya tiba, \qty{86}{\percent}; pada $X = 2.31$, $e^{-0.147\times2.31} =
e^{-0.340} = 0.712$, \qty{71}{\percent}. Pada $z = \ang{85}$, $\sec z =
11.5$; tetapi model sejajar-bidang membuat lintasan tak hingga di horizon,
sedangkan lintasan melalui atmosfer yang melengkung tetap berhingga, dan refraksi juga membelokkan sinar: di dekat horizon rumus itu melebih-lebihkan \omterm{def:b1:observing-the-sky:extinction}{massa udara} dan galatnya tumbuh dengan cepat.
\end{solution}

\begin{solution}{exo:b1:observing-the-sky:11}
(a) $\mathrm{SNR} = Rt/\sqrt{(R + B)t}$, sehingga $t = \mathrm{SNR}^2(R + B)/R^2
= 100\times2050/2500 = \qty{82}{s}$. (b) Membagi dua \omterm{def:b1:observing-the-sky:seeing}{kekaburan} membagi
luas citra, dan karenanya $B$, dengan 4: $B = 500$ dan $t = 100\times550/2500 =
\qty{22}{s}$. (c) Untuk $B \gg R$, $t \approx \mathrm{SNR}^2B/R^2$. Laju sumber
$R$ tumbuh sebagai $D^2$; laju langit $B$ tumbuh sebagai luas pengumpulan
$D^2$ dikali luas langit citra, $\theta^2$. Maka $t \propto
\theta^2D^2/D^4 = \theta^2/D^2$: membagi dua \omterm{def:b1:observing-the-sky:seeing}{kekaburan} sama bernilainya dengan menggandakan diameter.
\end{solution}

\begin{solution}{exo:b1:observing-the-sky:12}
(a) $N = 1200\times100 = \num{120000}$ garis; pada orde pertama daya resolusinya
\num{1.2e5}, dan $\Delta\lambda = 550/(\num{1.2e5}) =
\qty{4.6e-3}{nm}$. (b) $v = c\,\Delta\lambda/\lambda = c/(\num{1.2e5}) =
\qty{2.5}{km.s^{-1}}$. (c) $\Delta\lambda = 550\times29.78/\num{299792} =
\qty{0.055}{nm}$, sekitar 12 elemen resolusi. Pengamat harus menghitung, untuk tanggal dan waktu pengamatan, komponen kecepatan orbit Bumi ke arah bintang, lalu mengurangkannya, untuk memperoleh kecepatan bintang relatif terhadap Matahari.
\end{solution}

\begin{solution}{pb:b1:observing-the-sky:1}
\textbf{1.} $d = 1/0.37921 = \qty{2.637}{\parsec}$.

\textbf{2.} $2.637\times206265 = \qty{5.44e5}{\au}$;
$2.637\times\num{3.0857e16} = \qty{8.14e16}{m}$.

\textbf{3.} $2.637\times3.2616 = \qty{8.60}{\lightyear}$: cahaya yang kita lihat malam ini meninggalkan Sirius 8.6 tahun yang lalu.

\textbf{4.} Arah itu berayun sebesar $2p = \ang{;;0.758}$, sekitar dua kali \omterm{def:b1:observing-the-sky:seeing}{kekaburan}. Ayunan diukur dari posisi pusat citra, bukan dari tepinya: dirata-ratakan atas banyak foton, pusat itu dapat ditentukan hingga pecahan kecil dari lebar citra, makin teliti semakin banyak foton yang dihitung.

\textbf{5.} $\Delta d/d = \Delta p/p = 0.001/0.37921 = \qty{0.26}{\percent}$,
yaitu $8.60\times0.0026 = \qty{0.023}{\lightyear}$.

\textbf{6.} $10^{0.4\times(0.03 + 1.46)} = 10^{0.596} = 3.94$.

\textbf{7.} $10^{0.4\times(26.74 - 1.46)} = 10^{10.112} = \num{1.29e10}$.

\textbf{8.} $10^{0.4\times(8.44 + 1.46)} = 10^{3.96} = \num{9120}$.

\textbf{9.} $z = \ang{90} - \ang{82.09} = \ang{7.91}$, $X = \sec\ang{7.91} =
1.0096$; ekstingsi $0.15\times1.0096 = 0.151$ magnitudo; diukur
$m = -1.46 + 0.151 = -1.31$.

\textbf{10.} $z = \ang{60}$, $X = 2.000$; ekstingsi 0.30 magnitudo,
$m = -1.16$. Pecahan yang diteruskan adalah $10^{-0.4\times0.30} = 0.759$:
\qty{24}{\percent} cahaya hilang.

\textbf{11.} Saat terbit $z = \ang{90}$ dan $\sec z$ tak hingga, sedangkan
lintasan sebenarnya melalui atmosfer yang melengkung panjang tetapi berhingga; refraksi juga membelokkan sinar. Di dekat horizon ekstingsi besar dan berubah cepat, dan koreksi sejajar-bidang keliru.

\textbf{12.} $\num{3.0e4}\times10^{0.4\times(10.0 - 8.44)} = \num{3.0e4}
\times10^{0.624} = \num{3.0e4}\times4.21 = \num{1.26e5}$ fotoelektron per
detik.

\textbf{13.} SNR 1000 memerlukan $N = 1000^2 = \num{1.0e6}$
fotoelektron: $t = \num{1.0e6}/\num{1.26e5} = \qty{7.9}{s}$. Galat magnitudonya
adalah $1.086/1000 = \num{0.0011}$.

\textbf{14.} $\num{3.0e4}\times10^{0.4\times(10.0 + 1.46)} = \num{3.0e4}
\times10^{4.584} = \num{1.15e9}$ fotoelektron per detik.

\textbf{15.} $0.20\times\num{1.15e9} = \num{2.30e8}$ per detik di piksel terterang, yang terisi dalam $\num{1.0e5}/\num{2.30e8} =
\qty{4.3e-4}{s}$, kurang dari setengah milidetik.

\textbf{16.} Sayap Sirius memberikan $\num{1.15e9}\times10^{-4} =
\num{1.15e5}$ per detik; Sirius~B memberikan $0.20\times\num{1.26e5} =
\num{2.52e4}$ per detik. Silau Sirius 4.6 kali lebih terang daripada puncak Sirius~B.

\textbf{17.} Piksel itu menerima $\num{1.15e5} + \num{2.52e4} = \num{1.40e5}$
per detik dan terisi dalam $\num{1.0e5}/\num{1.40e5} = \qty{0.71}{s}$. Dalam
\qty{0.5}{s}: sinyal $\num{2.52e4}\times0.5 = \num{1.26e4}$, total
$\num{7.01e4}$, derau $\sqrt{\num{7.01e4}} = 265$, sehingga $\mathrm{SNR} =
\num{1.26e4}/265 = 48$. Sirius~B dapat dideteksi sejauh menyangkut derau foton; yang membuatnya sulit adalah silau Sirius harus diketahui dan dikurangkan dengan ketelitian jauh lebih baik daripada seperlima nilainya, sementara inti Sirius jenuh dalam kurang dari satu milidetik.

\textbf{18.} $m - M = 5\log_{10}(2.637/10) = 5\times(-0.579) = -2.89$.

\textbf{19.} $M = -1.46 + 2.89 = 1.43$.

\textbf{20.} $d = 1/0.13023 = \qty{7.679}{\parsec}$;
$m - M = 5\log_{10}0.7679 = -0.57$, sehingga $M = 0.03 + 0.57 = 0.60$.

\textbf{21.} Vega, dengan \omterm{def:b1:observing-the-sky:absolute}{magnitudo mutlak} yang lebih kecil, dengan faktor
$10^{0.4\times(1.43 - 0.60)} = 10^{0.332} = 2.15$. Sirius tampak lebih terang hanya karena ia lebih dekat.

\textbf{22.} $m = M_\odot + (m - M) = 4.83 - 2.89 = 1.94$: Matahari akan menjadi bintang yang cukup terang untuk dilihat dengan mata telanjang.

\textbf{23.} $5\log_{10}(d/\qty{10}{\parsec}) = 6.0 - 1.43 = 4.57$, sehingga
$d = 10\times10^{0.914} = \qty{82.0}{\parsec}$, sekitar \qty{267}{\lightyear}.

\textbf{24.} $L/L_\odot = 10^{0.4\times(4.83 - 1.43)} = 10^{1.36} = 22.9$:
dalam pita V, Sirius melampaui \textbf{sekitar 23 matahari}.
\end{solution}
